%========================================
% MatinBook - Stage 09: Algorithm Test
% Test: Algorithm Environments + Captions + List of Algorithms + References
% Purpose: Validate algorithm display, numbering, and indexing
%========================================

% Add parent directory to search path
\makeatletter
\def\input@path{{../}}
\makeatother

\documentclass{matinbook}

\title{آزمون سیستم الگوریتم‌ها}
\author{تیم توسعه MatinBook}
\date{تابستان ۱۴۰۵}

\begin{document}

%----------------------------------------
% Title Page
%----------------------------------------
\maketitle

%----------------------------------------
% List of Algorithms
%----------------------------------------
\listofalgorithms
\clearpage

%----------------------------------------
% Chapter 1: Search Algorithms
%----------------------------------------
\chapter{الگوریتم‌های جستجو}

\section{مقدمه}

الگوریتم‌های جستجو\index{الگوریتم!جستجو} از پایه‌ای‌ترین
الگوریتم‌های علوم کامپیوتر هستند. در این فصل،
دو الگوریتم اصلی جستجو را بررسی می‌کنیم.

\section{جستجوی خطی}

ساده‌ترین روش جستجو، بررسی تک‌تک عناصر آرایه است.
این روش برای آرایه‌های مرتب و نامرتب کار می‌کند:

\begin{latin}
\begin{algorithm}
\begin{algorithmic}[1]
    \Require $A[1..n]$: array of $n$ elements
    \Require $x$: target value to find
    \Ensure Index $i$ such that $A[i] = x$, or $-1$ if not found
    
    \For{$i \gets 1$ \textbf{to} $n$}
        \If{$A[i] = x$}
            \State \Return $i$
        \EndIf
    \EndFor
    \State \Return $-1$
\end{algorithmic}
\end{algorithm}
\algcaption[alg:linear-search]{جستجوی خطی}
\end{latin}

\begin{theorem}[پیچیدگی جستجوی خطی]
    الگوریتم \algref{alg:linear-search} در بدترین حالت
    $O(n)$ مقایسه انجام می‌دهد، زیرا ممکن است مجبور شود
    تمام $n$ عنصر آرایه را بررسی کند.
    \label{thm:linear-complexity}
\end{theorem}

\section{جستجوی دودویی}

برای آرایه‌های \textbf{مرتب}، جستجوی دودویی
بسیار کارآمدتر از جستجوی خطی است:

\begin{latin}
\begin{algorithm}
\begin{algorithmic}[1]
    \Require $A[1..n]$: sorted array (ascending order)
    \Require $x$: target value to find
    \Ensure Index $i$ such that $A[i] = x$, or $-1$ if not found
    
    \State $left \gets 1$
    \State $right \gets n$
    
    \While{$left \leq right$}
        \State $mid \gets \lfloor (left + right) / 2 \rfloor$
        \If{$A[mid] = x$}
            \State \Return $mid$
        \ElsIf{$A[mid] < x$}
            \State $left \gets mid + 1$
        \Else
            \State $right \gets mid - 1$
        \EndIf
    \EndWhile
    
    \State \Return $-1$
\end{algorithmic}
\end{algorithm}
\algcaption[alg:binary-search]{جستجوی دودویی}
\end{latin}

\begin{theorem}[پیچیدگی جستجوی دودویی]
    الگوریتم \algref{alg:binary-search} در بدترین حالت
    $O(\log n)$ مقایسه انجام می‌دهد، زیرا در هر مرحله
    فضای جستجو را نصف می‌کند.
    \label{thm:binary-complexity}
\end{theorem}

\section{مقایسه دو الگوریتم}

مقایسه \algref{alg:linear-search} و \algref{alg:binary-search}
نشان می‌دهد که برای آرایه‌های مرتب، جستجوی دودویی
به مراتب سریع‌تر است:

\begin{table}[h]
\centering
\caption{مقایسه الگوریتم‌های جستجو}
\label{tab:search-comparison}
\begin{tabular}{|c|C{3cm}|C{3cm}|c|}
\hline
\textbf{معیار} & \textbf{جستجوی خطی} & \textbf{جستجوی دودویی} & \textbf{برنده} \\
\hline
پیچیدگی زمانی & $O(n)$ & $O(\log n)$ & دودویی \\
نیاز به مرتب‌سازی & خیر & \textcolor{matinred}{بله} & خطی \\
پیاده‌سازی & ساده & متوسط & خطی \\
کارایی برای $n$ کوچک & خوب & سربار بیشتر & خطی \\
کارایی برای $n$ بزرگ & ضعیف & \textcolor{matingreen}{عالی} & دودویی \\
\hline
\end{tabular}
\end{table}

%----------------------------------------
% Chapter 2: Sorting Algorithms
%----------------------------------------
\chapter{الگوریتم‌های مرتب‌سازی}

\section{مرتب‌سازی سریع}

\lr{Quick Sort}\index{الگوریتم!مرتب‌سازی سریع}
یکی از سریع‌ترین الگوریتم‌های مرتب‌سازی در عمل است:

\begin{latin}
\begin{algorithm}
\begin{algorithmic}[1]
    \Require $A[1..n]$: array of $n$ elements
    \Require $low, high$: indices defining the subarray
    \Ensure Array $A$ is sorted in ascending order
    
    \Function{QuickSort}{$A$, $low$, $high$}
        \If{$low < high$}
            \State $p \gets \Call{Partition}{A, low, high}$
            \State $\Call{QuickSort}{A, low, p - 1}$
            \State $\Call{QuickSort}{A, p + 1, high}$
        \EndIf
    \EndFunction
    \Statex
    \Function{Partition}{$A$, $low$, $high$}
        \State $pivot \gets A[high]$
        \State $i \gets low - 1$
        \For{$j \gets low$ \textbf{to} $high - 1$}
            \If{$A[j] \leq pivot$}
                \State $i \gets i + 1$
                \State Swap $A[i]$ and $A[j]$
            \EndIf
        \EndFor
        \State Swap $A[i+1]$ and $A[high]$
        \State \Return $i + 1$
    \EndFunction
\end{algorithmic}
\end{algorithm}
\algcaption[alg:quicksort]{مرتب‌سازی سریع (\lr{Quick Sort})}
\end{latin}

\section{الگوریتم اقلیدس}

الگوریتم اقلیدس\index{الگوریتم!اقلیدس}
برای محاسبه بزرگترین مقسوم‌علیه مشترک (\lr{GCD}):

\begin{latin}
\begin{algorithm}
\begin{algorithmic}[1]
    \Require $a, b$: non-negative integers
    \Ensure $\gcd(a, b)$: greatest common divisor
    
    \While{$b \neq 0$}
        \State $r \gets a \bmod b$
        \State $a \gets b$
        \State $b \gets r$
    \EndWhile
    \State \Return $a$
\end{algorithmic}
\end{algorithm}
\algcaption[alg:euclidean]{الگوریتم اقلیدس برای $\gcd(a, b)$}
\end{latin}

\begin{example}
    محاسبه $\gcd(48, 18)$ با الگوریتم \algref{alg:euclidean}:
    \begin{align*}
        \gcd(48, 18) &: \; 48 \bmod 18 = 12 \\
        \gcd(18, 12) &: \; 18 \bmod 12 = 6 \\
        \gcd(12, 6)  &: \; 12 \bmod 6 = 0 \\
        \gcd(6, 0)   &: \; \text{نتیجه} = 6
    \end{align*}
    \label{ex:gcd-example}
\end{example}

%----------------------------------------
% Chapter 3: Graph Algorithms
%----------------------------------------
\chapter{الگوریتم‌های گراف}

\section{جستجوی اول سطح (\lr{BFS})}

\lr{Breadth-First Search}\index{الگوریتم!BFS}
گراف را سطح به سطح پیمایش می‌کند:

\begin{latin}
\begin{algorithm}
\begin{algorithmic}[1]
    \Require $G = (V, E)$: an unweighted graph
    \Require $s$: source vertex
    \Ensure $dist[v]$: shortest distance from $s$ to each $v \in V$
    
    \ForAll{$v \in V$}
        \State $dist[v] \gets \infty$
        \State $visited[v] \gets \text{false}$
    \EndFor
    
    \State $dist[s] \gets 0$
    \State $visited[s] \gets \text{true}$
    \State Create an empty queue $Q$
    \State $Q.\text{enqueue}(s)$
    
    \While{$Q$ is not empty}
        \State $u \gets Q.\text{dequeue}()$
        \ForAll{neighbor $v$ of $u$}
            \If{$\neg visited[v]$}
                \State $visited[v] \gets \text{true}$
                \State $dist[v] \gets dist[u] + 1$
                \State $Q.\text{enqueue}(v)$
            \EndIf
        \EndFor
    \EndWhile
    
    \State \Return $dist$
\end{algorithmic}
\end{algorithm}
\algcaption[alg:bfs]{جستجوی اول سطح (\lr{BFS})}
\end{latin}

\section{الگوریتم دایکسترا}

الگوریتم دایکسترا\index{الگوریتم!دایکسترا}
کوتاه‌ترین مسیر در گراف وزن‌دار را پیدا می‌کند:

\begin{latin}
\begin{algorithm}
\begin{algorithmic}[1]
    \Require $G = (V, E)$: weighted graph with non-negative weights
    \Require $s$: source vertex
    \Ensure $dist[v]$: shortest distance from $s$ to each $v \in V$
    
    \ForAll{$v \in V$}
        \State $dist[v] \gets \infty$
        \State $prev[v] \gets \text{null}$
    \EndFor
    \State $dist[s] \gets 0$
    \State $PQ \gets$ priority queue of vertices keyed by $dist$
    \State Insert all vertices into $PQ$
    
    \While{$PQ$ is not empty}
        \State $u \gets$ extract vertex with minimum $dist$ from $PQ$
        \ForAll{neighbor $v$ of $u$}
            \State $alt \gets dist[u] + weight(u, v)$
            \If{$alt < dist[v]$}
                \State $dist[v] \gets alt$
                \State $prev[v] \gets u$
                \State Decrease key of $v$ in $PQ$ to $alt$
            \EndIf
        \EndFor
    \EndWhile
    
    \State \Return $(dist, prev)$
\end{algorithmic}
\end{algorithm}
\algcaption[alg:dijkstra]{الگوریتم دایکسترا (\lr{Dijkstra's Algorithm})}
\end{latin}

%----------------------------------------
% Chapter 4: Algorithm Analysis
%----------------------------------------
\chapter{تحلیل و مقایسه الگوریتم‌ها}

\section{ارجاع به الگوریتم‌ها}

سیستم ارجاع \lr{MatinBook} امکان ارجاع به الگوریتم‌ها
را با استفاده از \texttt{\textbackslash algref} فراهم می‌کند:

\begin{itemize}
    \item \algref{alg:linear-search} --- ساده‌ترین روش جستجو
    با پیچیدگی $O(n)$
    \item \algref{alg:binary-search} --- جستجوی سریع
    با پیچیدگی $O(\log n)$ برای آرایه‌های مرتب
    \item \algref{alg:quicksort} --- مرتب‌سازی سریع
    با پیچیدگی متوسط $O(n \log n)$
    \item \algref{alg:euclidean} --- الگوریتم باستانی اقلیدس
    با پیچیدگی $O(\log \min(a,b))$
    \item \algref{alg:bfs} --- پیمایش سطح به سطح گراف
    با پیچیدگی $O(V + E)$
    \item \algref{alg:dijkstra} --- کوتاه‌ترین مسیر
    با پیچیدگی $O((V+E) \log V)$
\end{itemize}

\section{جدول مقایسه}

\begin{table}[h]
\centering
\caption{مقایسه پیچیدگی الگوریتم‌ها}
\label{tab:algorithm-comparison}
\begin{tabular}{|c|C{3cm}|c|c|}
\hline
\textbf{الگوریتم} & \textbf{کاربرد} & \textbf{پیچیدگی} & \textbf{نیازمندی} \\
\hline
جستجوی خطی & جستجوی ساده & $O(n)$ & ندارد \\
جستجوی دودویی & جستجوی سریع & $O(\log n)$ & آرایه مرتب \\
مرتب‌سازی سریع & مرتب‌سازی & $O(n \log n)$ & ندارد \\
الگوریتم اقلیدس & محاسبه GCD & $O(\log n)$ & اعداد صحیح \\
\lr{BFS} & پیمایش گراف & $O(V+E)$ & گراف بدون وزن \\
دایکسترا & کوتاه‌ترین مسیر & $O((V+E)\log V)$ & یال‌های غیرمنفی \\
\hline
\end{tabular}
\end{table}

%----------------------------------------
% Summary
%----------------------------------------
\chapter{خلاصه نتایج}

\section{وضعیت سیستم الگوریتم}

\begin{table}[h]
\centering
\caption{وضعیت قابلیت‌های الگوریتم}
\label{tab:algorithm-status}
\begin{tabular}{|c|C{6cm}|c|}
\hline
\textbf{قابلیت} & \textbf{توضیحات} & \textbf{وضعیت} \\
\hline
شماره‌گذاری & شماره‌گذاری خودکار به صورت «فصل.شماره» & \textcolor{matingreen}{✅} \\
عنوان فارسی & نمایش عنوان با \lr{\textbackslash algcaption} & \textcolor{matingreen}{✅} \\
فهرست الگوریتم‌ها & \lr{\textbackslash listofalgorithms} & \textcolor{matingreen}{✅} \\
ارجاع & \lr{\textbackslash algref} با پیشوند «الگوریتم» & \textcolor{matingreen}{✅} \\
\lr{Require/Ensure} & پیش‌نیاز و خروجی الگوریتم & \textcolor{matingreen}{✅} \\
\lr{For/While/If} & ساختارهای کنترلی استاندارد & \textcolor{matingreen}{✅} \\
\lr{Function} & تعریف و فراخوانی تابع & \textcolor{matingreen}{✅} \\
محیط \lr{latin} & شبه‌کد درون محیط لاتین & \textcolor{matingreen}{✅} \\
\hline
\end{tabular}
\end{table}

\section{نتیجه}

\textbf{ماژول الگوریتم \lr{MatinBook} نسخه ۱ با موفقیت آزمایش شد.}
۶ الگوریتم مختلف با عناوین فارسی، شماره‌گذاری خودکار،
فهرست الگوریتم‌ها و ارجاع هوشمند به درستی کار می‌کنند. 🎉

\end{document}